% Beweise zu den unverändert im Hauptband deklarierten Aussagen.
\section{Koordinatenträger und Koordinatenprodukt}

Diese Ergänzung beweist alle sechzehn Sätze des Abschnitts
\emph{\(e\)-Koordinaten\-abbildung} in Band~31.
Die verlinkten Satznummern gehören zum Hauptband. Für jede Tabelle gelten
die folgenden bloßen Schreibabkürzungen mit denselben Parametern
\(A,\star,e\):
\[
\begin{aligned}
L&=\{p\in A\mid\exists x\in A\;(p=x\star e)\},\\
R&=\{q\in A\mid\exists y\in A\;(q=e\star y)\},\\
K&=\SemigroupCoordinateCarrier{A}{\star}{e},\\
\odot&=\SemigroupCoordinateProductOp{A}{\star}{e},\\
\chi&=\SemigroupCoordinateMap{A}{\star}{e},\qquad
\mu=\SemigroupCoordinateMultiplicationMap{A}{\star}{e}.
\end{aligned}
\]
Die Gleichung \(K=L\times R\) folgt unter den Voraussetzungen
\(\SemigroupAlg{A}{\star}\) und \(e\in A\) aus
\FormulaRefAuto{SemigroupCoordinateCarrierDef}[def].
Kommagetrennte Mitgliedschaften bedeuten rechtsassoziierte Konjunktionen.
Die Komponenten einer Produktzerlegung werden als frische Eigenvariablen
eingeführt und mit \(\exists E^{*}\) wieder entlassen. Die dabei benötigten
Konjunktionsglieder werden vor weiteren Regelanwendungen einzeln gewonnen.
Bei \(=E\) stehen zunächst die gerichteten Gleichheiten und zuletzt die
Ausgangsformel; eine erforderliche Umkehrung wird eigens begründet.
Die Definitionen der drei Operationen beziehungsweise Abbildungen enthalten
bereits ihre Typisierung. Die folgenden Tabellen überprüfen zusätzlich
die jeweils verwendeten Zielbereiche.

\subsection{Elementkriterium}
\hypertarget{coordinates.proof.SemigroupCoordinateCarrierMembership}{}
\noindent Beweis zu
\FormulaRefAuto{SemigroupCoordinateCarrierMembership}[thm].
\begin{tabproofwide}
  \proofstepwidestar[1]{\SemigroupAlg{A}{\star}}{\rA}
  \proofstepwidestar[2]{e\in A}{\rA}
  \proofstepwidestar[1,2]{K=L\times R}{\FormulaRefAuto{SemigroupCoordinateCarrierDef}[def]{1,2}}
  \proofstepwidestar[1,2]{L\times R=K}{\FormulaRefAuto{a=b\vdash b=a}[thm]{3}}
  \proofstepwidestar[5]{(p,q)\in K}{\rA}
  \proofstepwidestar[1,2,5]{(p,q)\in L\times R}{\rIE{3,5}}
  \proofstepwidestar[1,2,5]{p\in L\land q\in R}{\FormulaRefAuto{(a,b)\in A\times B\eqvdash a\in A\land b\in B}[thm]{6}}
  \proofstepwidestar[1,2,5]{p\in L}{\rAEa{7}}
  \proofstepwidestar[1,2,5]{q\in R}{\rAEb{7}}
  \proofstepwidestar[1,2,5]{p\in A\land\exists x\in A\;(p=x\star e)}{\FormulaRefAuto{x\in\{u\in A\mid P(u)\}\eqvdash x\in A\land P(x)}[thm]{8}}
  \proofstepwidestar[1,2,5]{q\in A\land\exists y\in A\;(q=e\star y)}{\FormulaRefAuto{x\in\{u\in A\mid P(u)\}\eqvdash x\in A\land P(x)}[thm]{9}}
  \proofstepwidestar[1,2,5]{\begin{aligned}[t]&(p\in A\land\exists x\in A\;(p=x\star e))\\[-2pt]&\quad\land(q\in A\land\exists y\in A\;(q=e\star y))\end{aligned}}{\rAI{10,11}}
  \proofstepwidestar[1,2]{\begin{aligned}[t](p,q)\in K\rightarrow{}&(p\in A\land\exists x\in A\;(p=x\star e))\\[-2pt]&\land(q\in A\land\exists y\in A\;(q=e\star y))\end{aligned}}{\rRI{5,12}}
  \proofstepwidestar[14]{\begin{aligned}[t]&(p\in A\land\exists x\in A\;(p=x\star e))\\[-2pt]&\quad\land(q\in A\land\exists y\in A\;(q=e\star y))\end{aligned}}{\rA}
  \proofstepwidestar[14]{p\in A\land\exists x\in A\;(p=x\star e)}{\rAEa{14}}
  \proofstepwidestar[14]{q\in A\land\exists y\in A\;(q=e\star y)}{\rAEb{14}}
  \proofstepwidestar[14]{p\in L}{\FormulaRefAuto{x\in\{u\in A\mid P(u)\}\eqvdash x\in A\land P(x)}[thm]{15}}
  \proofstepwidestar[14]{q\in R}{\FormulaRefAuto{x\in\{u\in A\mid P(u)\}\eqvdash x\in A\land P(x)}[thm]{16}}
  \proofstepwidestar[14]{p\in L\land q\in R}{\rAI{17,18}}
  \proofstepwidestar[14]{(p,q)\in L\times R}{\FormulaRefAuto{(a,b)\in A\times B\eqvdash a\in A\land b\in B}[thm]{19}}
  \proofstepwidestar[1,2,14]{(p,q)\in K}{\rIE{4,20}}
  \proofstepwidestar[1,2]{\begin{aligned}[t]&(p\in A\land\exists x\in A\;(p=x\star e))\\[-2pt]&\quad\land(q\in A\land\exists y\in A\;(q=e\star y))\rightarrow(p,q)\in K\end{aligned}}{\rRI{14,21}}
  \proofstepwidestar[1,2]{\begin{aligned}[t](p,q)\in K\leftrightarrow{}&(p\in A\land\exists x\in A\;(p=x\star e))\\[-2pt]&\land(q\in A\land\exists y\in A\;(q=e\star y))\end{aligned}}{\rLRI{13,22}}
\end{tabproofwide}

\Needspace{22\baselineskip}
\subsection{Nichtleerer Koordinatenträger}
\hypertarget{coordinates.proof.SemigroupCoordinateCarrierNonempty}{}
\noindent Beweis zu
\FormulaRefAuto{SemigroupCoordinateCarrierNonempty}[thm].
\begin{tabproofwide}
  \proofstepwidestar[1]{\SemigroupAlg{A}{\star}}{\rA}
  \proofstepwidestar[2]{e\in A}{\rA}
  \proofstepwidestar[1,2]{e\star e\in A}{\FormulaRefAuto{SemigroupClosureAxiom}[ax]{1,\rAI{2,2}}}
  \proofstepwidestar[2]{\exists x\in A\;(e\star e=x\star e)}{\rEI{\rAI{2,\rII}}}
  \proofstepwidestar[2]{\exists y\in A\;(e\star e=e\star y)}{\rEI{\rAI{2,\rII}}}
  \proofstepwidestar[1,2]{\begin{aligned}[t]&(e\star e,e\star e)\in K\leftrightarrow{}\\[-2pt]&\quad(e\star e\in A\land\exists x\in A\;(e\star e=x\star e))\\[-2pt]&\quad\land(e\star e\in A\land\exists y\in A\;(e\star e=e\star y))\end{aligned}}{\FormulaRefAuto{SemigroupCoordinateCarrierMembership}[thm]{1,2}}
  \proofstepwidestar[1,2]{(e\star e,e\star e)\in K}{\rRE{\rLREb{6},\rAI{\rAI{3,4},\rAI{3,5}}}}
  \proofstepwidestar[1,2]{K\neq\varnothing}{\FormulaRefAuto{a\in A\vdash A\neq\varnothing}[thm]{7}}
\end{tabproofwide}

\subsection{Auswertung des Produkts}
\hypertarget{coordinates.proof.SemigroupCoordinateProductEvaluation}{}
\noindent Beweis zu
\FormulaRefAuto{SemigroupCoordinateProductEvaluation}[thm].
\begin{tabproofwide}
  \proofstepwidestar[1]{\SemigroupAlg{A}{\star}}{\rA}
  \proofstepwidestar[2]{e\in A}{\rA}
  \proofstepwidestar[3]{(p,q),(r,s)\in K}{\rA}
  \proofstepwidestar[1,2,3]{(p,q)\odot(r,s)=(p,s)}{\FormulaRefAuto{SemigroupCoordinateProductDef}[def]{1,2,3}}
\end{tabproofwide}

\Needspace{14\baselineskip}
\subsection{Halbgruppeneigenschaft des Koordinatenprodukts}
\hypertarget{coordinates.proof.SemigroupCoordinateProductSemigroup}{}
\noindent Beweis zu
\FormulaRefAuto{SemigroupCoordinateProductSemigroup}[thm].
\begin{tabproofsplitwide}
\proofpartwide{Abgeschlossenheit mit vollständiger Produktzerlegung}
  \proofstepwidestar[1]{\SemigroupAlg{A}{\star}}{\rA}
  \proofstepwidestar[2]{e\in A}{\rA}
  \proofstepwidestar[1,2]{K=L\times R}{\FormulaRefAuto{SemigroupCoordinateCarrierDef}[def]{1,2}}
  \proofstepwidestar[1,2]{L\times R=K}{\FormulaRefAuto{a=b\vdash b=a}[thm]{3}}
  \proofstepwidestar[1,2]{\odot\colon K\times K\to K}{\FormulaRefAuto{SemigroupCoordinateProductDef}[def]{1,2}}
  \proofstepwidestar[6]{u\in K}{\rA}
  \proofstepwidestar[7]{v\in K}{\rA}
  \proofstepwidestar[1,2,6]{u\in L\times R}{\rIE{3,6}}
  \proofstepwidestar[1,2,6]{\exists p\in L\,\exists q\in R\;(u=(p,q))}{\FormulaRefAuto{x\in A\times B\eqvdash\exists a\in A\,\exists b\in B\,x=(a,b)}[thm]{8}}
  \proofstepwidestar[10]{p\in L\land(q\in R\land u=(p,q))}{\rA}
  \proofstepwidestar[10]{p\in L}{\rAEa{10}}
  \proofstepwidestar[10]{q\in R}{\rAEa{\rAEb{10}}}
  \proofstepwidestar[10]{u=(p,q)}{\rAEb{\rAEb{10}}}
  \proofstepwidestar[6,10]{(p,q)\in K}{\rIE{13,6}}
  \proofstepwidestar[1,2,7]{v\in L\times R}{\rIE{3,7}}
  \proofstepwidestar[1,2,7]{\exists r\in L\,\exists s\in R\;(v=(r,s))}{\FormulaRefAuto{x\in A\times B\eqvdash\exists a\in A\,\exists b\in B\,x=(a,b)}[thm]{15}}
  \proofstepwidestar[17]{r\in L\land(s\in R\land v=(r,s))}{\rA}
  \proofstepwidestar[17]{r\in L}{\rAEa{17}}
  \proofstepwidestar[17]{s\in R}{\rAEa{\rAEb{17}}}
  \proofstepwidestar[17]{v=(r,s)}{\rAEb{\rAEb{17}}}
  \proofstepwidestar[7,17]{(r,s)\in K}{\rIE{20,7}}
  \proofstepwidestar[10,17]{(p,s)\in L\times R}{\FormulaRefAuto{(a,b)\in A\times B\eqvdash a\in A\land b\in B}[thm]{\rAI{11,19}}}
  \proofstepwidestar[1,2,10,17]{(p,s)\in K}{\rIE{4,22}}
  \proofstepwidestar[1,2,6,7,10,17]{(p,q)\odot(r,s)=(p,s)}{\FormulaRefAuto{SemigroupCoordinateProductEvaluation}[thm]{1,2,\rAI{14,21}}}
  \proofstepwidestar[10,17]{u\odot v=(p,q)\odot(r,s)}{\rIE{13,20,\rII}}
  \proofstepwidestar[1,2,6,7,10,17]{u\odot v=(p,s)}{\rChain{25,24}}
  \proofstepwidestar[1,2,6,7,10,17]{(p,s)=u\odot v}{\FormulaRefAuto{a=b\vdash b=a}[thm]{26}}
  \proofstepwidestar[1,2,6,7,10,17]{u\odot v\in K}{\rIE{27,23}}
  \proofstepwidestar[1,2,6,7,10]{u\odot v\in K}{\rEEStar{16,17,28}}
  \proofstepwidestar[1,2,6,7]{u\odot v\in K}{\rEEStar{9,10,29}}
  \proofstepwidestar[1,2,6]{\forall v\in K\;(u\odot v\in K)}{\rUI{\rRI{7,30}}}
\nopagebreak
  \proofstepwidestar[1,2]{\forall u,v\in K\;(u\odot v\in K)}{\rUI{\rRI{6,31}}}
\closeproofpartwide
\Needspace{8\baselineskip}
\proofpartwide{Assoziativität mit allen Zwischenprodukten}
\setcounter{proofstepnr}{32}
  \proofstepwidestar[33]{x\in K}{\rA}
  \proofstepwidestar[34]{y\in K}{\rA}
  \proofstepwidestar[35]{z\in K}{\rA}
  \proofstepwidestar[1,2,33]{x\in L\times R}{\rIE{3,33}}
  \proofstepwidestar[1,2,33]{\exists a\in L\,\exists b\in R\;(x=(a,b))}{\FormulaRefAuto{x\in A\times B\eqvdash\exists a\in A\,\exists b\in B\,x=(a,b)}[thm]{36}}
  \proofstepwidestar[38]{a\in L\land(b\in R\land x=(a,b))}{\rA}
  \proofstepwidestar[38]{a\in L}{\rAEa{38}}
  \proofstepwidestar[38]{b\in R}{\rAEa{\rAEb{38}}}
  \proofstepwidestar[38]{x=(a,b)}{\rAEb{\rAEb{38}}}
  \proofstepwidestar[33,38]{(a,b)\in K}{\rIE{41,33}}
  \proofstepwidestar[1,2,34]{y\in L\times R}{\rIE{3,34}}
  \proofstepwidestar[1,2,34]{\exists c\in L\,\exists d\in R\;(y=(c,d))}{\FormulaRefAuto{x\in A\times B\eqvdash\exists a\in A\,\exists b\in B\,x=(a,b)}[thm]{43}}
  \proofstepwidestar[45]{c\in L\land(d\in R\land y=(c,d))}{\rA}
  \proofstepwidestar[45]{c\in L}{\rAEa{45}}
  \proofstepwidestar[45]{d\in R}{\rAEa{\rAEb{45}}}
  \proofstepwidestar[45]{y=(c,d)}{\rAEb{\rAEb{45}}}
  \proofstepwidestar[34,45]{(c,d)\in K}{\rIE{48,34}}
  \proofstepwidestar[1,2,35]{z\in L\times R}{\rIE{3,35}}
  \proofstepwidestar[1,2,35]{\exists f\in L\,\exists g\in R\;(z=(f,g))}{\FormulaRefAuto{x\in A\times B\eqvdash\exists a\in A\,\exists b\in B\,x=(a,b)}[thm]{50}}
  \proofstepwidestar[52]{f\in L\land(g\in R\land z=(f,g))}{\rA}
  \proofstepwidestar[52]{f\in L}{\rAEa{52}}
  \proofstepwidestar[52]{g\in R}{\rAEa{\rAEb{52}}}
  \proofstepwidestar[52]{z=(f,g)}{\rAEb{\rAEb{52}}}
  \proofstepwidestar[35,52]{(f,g)\in K}{\rIE{55,35}}
  \proofstepwidestar[38,45]{(a,d)\in L\times R}{\FormulaRefAuto{(a,b)\in A\times B\eqvdash a\in A\land b\in B}[thm]{\rAI{39,47}}}
  \proofstepwidestar[1,2,38,45]{(a,d)\in K}{\rIE{4,57}}
  \proofstepwidestar[45,52]{(c,g)\in L\times R}{\FormulaRefAuto{(a,b)\in A\times B\eqvdash a\in A\land b\in B}[thm]{\rAI{46,54}}}
  \proofstepwidestar[1,2,45,52]{(c,g)\in K}{\rIE{4,59}}
  \proofstepwidestar[1,2,33,34,38,45]{(a,b)\odot(c,d)=(a,d)}{\FormulaRefAuto{SemigroupCoordinateProductEvaluation}[thm]{1,2,\rAI{42,49}}}
  \proofstepwidestar[1,2,34,35,45,52]{(c,d)\odot(f,g)=(c,g)}{\FormulaRefAuto{SemigroupCoordinateProductEvaluation}[thm]{1,2,\rAI{49,56}}}
  \proofstepwidestar[38,45,52]{(x\odot y)\odot z=((a,b)\odot(c,d))\odot(f,g)}{\rIE{41,48,55,\rII}}
  \proofstepwidestar[1,2,33,34,38,45,52]{(x\odot y)\odot z=(a,d)\odot(f,g)}{\rIE{61,63}}
  \proofstepwidestar[1,2,35,38,45,52]{(a,d)\odot(f,g)=(a,g)}{\FormulaRefAuto{SemigroupCoordinateProductEvaluation}[thm]{1,2,\rAI{58,56}}}
  \proofstepwidestar[1,2,33,34,35,38,45,52]{(x\odot y)\odot z=(a,g)}{\rChain{64,65}}
  \proofstepwidestar[38,45,52]{x\odot(y\odot z)=(a,b)\odot((c,d)\odot(f,g))}{\rIE{41,48,55,\rII}}
  \proofstepwidestar[1,2,34,35,38,45,52]{x\odot(y\odot z)=(a,b)\odot(c,g)}{\rIE{62,67}}
  \proofstepwidestar[1,2,33,38,45,52]{(a,b)\odot(c,g)=(a,g)}{\FormulaRefAuto{SemigroupCoordinateProductEvaluation}[thm]{1,2,\rAI{42,60}}}
  \proofstepwidestar[1,2,33,34,35,38,45,52]{x\odot(y\odot z)=(a,g)}{\rChain{68,69}}
  \proofstepwidestar[1,2,33,34,35,38,45,52]{(a,g)=x\odot(y\odot z)}{\FormulaRefAuto{a=b\vdash b=a}[thm]{70}}
  \proofstepwidestar[1,2,33,34,35,38,45,52]{(x\odot y)\odot z=x\odot(y\odot z)}{\rIE{71,66}}
  \proofstepwidestar[1,2,33,34,35,38,45]{(x\odot y)\odot z=x\odot(y\odot z)}{\rEEStar{51,52,72}}
  \proofstepwidestar[1,2,33,34,35,38]{(x\odot y)\odot z=x\odot(y\odot z)}{\rEEStar{44,45,73}}
  \proofstepwidestar[1,2,33,34,35]{(x\odot y)\odot z=x\odot(y\odot z)}{\rEEStar{37,38,74}}
  \proofstepwidestar[1,2,33,34]{\forall z\in K\;((x\odot y)\odot z=x\odot(y\odot z))}{\rUI{\rRI{35,75}}}
  \proofstepwidestar[1,2,33]{\forall y,z\in K\;((x\odot y)\odot z=x\odot(y\odot z))}{\rUI{\rRI{34,76}}}
  \proofstepwidestar[1,2]{\forall x,y,z\in K\;((x\odot y)\odot z=x\odot(y\odot z))}{\rUI{\rRI{33,77}}}
  \proofstepwidestar[1,2]{\SemigroupAlg{K}{\odot}}{\FormulaRefAuto{SemigroupDef}[def]{32,78}}
\closeproofpartwide
\end{tabproofsplitwide}

\Needspace{23\baselineskip}
\subsection{Idempotenz des Koordinatenprodukts}
\hypertarget{coordinates.proof.SemigroupCoordinateProductIdempotence}{}
\noindent Beweis zu
\FormulaRefAuto{SemigroupCoordinateProductIdempotence}[thm].
\begin{tabproofwide}
  \proofstepwidestar[1]{\SemigroupAlg{A}{\star}}{\rA}
  \proofstepwidestar[2]{e\in A}{\rA}
  \proofstepwidestar[1,2]{K=L\times R}{\FormulaRefAuto{SemigroupCoordinateCarrierDef}[def]{1,2}}
  \proofstepwidestar[1,2]{L\times R=K}{\FormulaRefAuto{a=b\vdash b=a}[thm]{3}}
  \proofstepwidestar[5]{u\in K}{\rA}
  \proofstepwidestar[1,2,5]{u\in L\times R}{\rIE{3,5}}
  \proofstepwidestar[1,2,5]{\exists p\in L\,\exists q\in R\;(u=(p,q))}{\FormulaRefAuto{x\in A\times B\eqvdash\exists a\in A\,\exists b\in B\,x=(a,b)}[thm]{6}}
  \proofstepwidestar[8]{p\in L\land(q\in R\land u=(p,q))}{\rA}
  \proofstepwidestar[8]{p\in L}{\rAEa{8}}
  \proofstepwidestar[8]{q\in R}{\rAEa{\rAEb{8}}}
  \proofstepwidestar[8]{u=(p,q)}{\rAEb{\rAEb{8}}}
  \proofstepwidestar[5,8]{(p,q)\in K}{\rIE{11,5}}
  \proofstepwidestar[1,2,5,8]{(p,q)\odot(p,q)=(p,q)}{\FormulaRefAuto{SemigroupCoordinateProductEvaluation}[thm]{1,2,\rAI{12,12}}}
  \proofstepwidestar[8]{u\odot u=(p,q)\odot(p,q)}{\rIE{11,\rII}}
  \proofstepwidestar[1,2,5,8]{u\odot u=(p,q)}{\rChain{14,13}}
  \proofstepwidestar[8]{(p,q)=u}{\FormulaRefAuto{a=b\vdash b=a}[thm]{11}}
  \proofstepwidestar[1,2,5,8]{u\odot u=u}{\rIE{16,15}}
  \proofstepwidestar[1,2,5]{u\odot u=u}{\rEEStar{7,8,17}}
\end{tabproofwide}

\subsection{Dreitermgesetz des Koordinatenprodukts}
\hypertarget{coordinates.proof.SemigroupCoordinateProductThreeTermLaw}{}
\noindent Beweis zu
\FormulaRefAuto{SemigroupCoordinateProductThreeTermLaw}[thm].
\begin{tabproofsplitwide}
\proofpartwide{Produktzerlegungen und Zwischenprodukte}
  \proofstepwidestar[1]{\SemigroupAlg{A}{\star}}{\rA}
  \proofstepwidestar[2]{e\in A}{\rA}
  \proofstepwidestar[1,2]{K=L\times R}{\FormulaRefAuto{SemigroupCoordinateCarrierDef}[def]{1,2}}
  \proofstepwidestar[1,2]{L\times R=K}{\FormulaRefAuto{a=b\vdash b=a}[thm]{3}}
  \proofstepwidestar[5]{u,v,w\in K}{\rA}
  \proofstepwidestar[5]{u\in K}{\rAEa{5}}
  \proofstepwidestar[5]{v\in K}{\rAEa{\rAEb{5}}}
  \proofstepwidestar[5]{w\in K}{\rAEb{\rAEb{5}}}
  \proofstepwidestar[1,2,5]{u\in L\times R}{\rIE{3,6}}
  \proofstepwidestar[1,2,5]{\exists p\in L\,\exists q\in R\;(u=(p,q))}{\FormulaRefAuto{x\in A\times B\eqvdash\exists a\in A\,\exists b\in B\,x=(a,b)}[thm]{9}}
  \proofstepwidestar[11]{p\in L\land(q\in R\land u=(p,q))}{\rA}
  \proofstepwidestar[11]{p\in L}{\rAEa{11}}
  \proofstepwidestar[11]{q\in R}{\rAEa{\rAEb{11}}}
  \proofstepwidestar[11]{u=(p,q)}{\rAEb{\rAEb{11}}}
  \proofstepwidestar[5,11]{(p,q)\in K}{\rIE{14,6}}
  \proofstepwidestar[1,2,5]{v\in L\times R}{\rIE{3,7}}
  \proofstepwidestar[1,2,5]{\exists r\in L\,\exists s\in R\;(v=(r,s))}{\FormulaRefAuto{x\in A\times B\eqvdash\exists a\in A\,\exists b\in B\,x=(a,b)}[thm]{16}}
  \proofstepwidestar[18]{r\in L\land(s\in R\land v=(r,s))}{\rA}
  \proofstepwidestar[18]{r\in L}{\rAEa{18}}
  \proofstepwidestar[18]{s\in R}{\rAEa{\rAEb{18}}}
  \proofstepwidestar[18]{v=(r,s)}{\rAEb{\rAEb{18}}}
  \proofstepwidestar[5,18]{(r,s)\in K}{\rIE{21,7}}
  \proofstepwidestar[1,2,5]{w\in L\times R}{\rIE{3,8}}
  \proofstepwidestar[1,2,5]{\exists t\in L\,\exists z\in R\;(w=(t,z))}{\FormulaRefAuto{x\in A\times B\eqvdash\exists a\in A\,\exists b\in B\,x=(a,b)}[thm]{23}}
  \proofstepwidestar[25]{t\in L\land(z\in R\land w=(t,z))}{\rA}
  \proofstepwidestar[25]{t\in L}{\rAEa{25}}
  \proofstepwidestar[25]{z\in R}{\rAEa{\rAEb{25}}}
  \proofstepwidestar[25]{w=(t,z)}{\rAEb{\rAEb{25}}}
  \proofstepwidestar[5,25]{(t,z)\in K}{\rIE{28,8}}
  \proofstepwidestar[11,18]{(p,s)\in L\times R}{\FormulaRefAuto{(a,b)\in A\times B\eqvdash a\in A\land b\in B}[thm]{\rAI{12,20}}}
  \proofstepwidestar[1,2,11,18]{(p,s)\in K}{\rIE{4,30}}
  \proofstepwidestar[1,2,5,11,18]{(p,q)\odot(r,s)=(p,s)}{\FormulaRefAuto{SemigroupCoordinateProductEvaluation}[thm]{1,2,\rAI{15,22}}}
\closeproofpartwide
\Needspace{23\baselineskip}
\proofpartwide{Vergleich der Produkte und Entlassung der Zeugen}
\setcounter{proofstepnr}{32}
  \proofstepwidestar[11,18,25]{(u\odot v)\odot w=((p,q)\odot(r,s))\odot(t,z)}{\rIE{14,21,28,\rII}}
  \proofstepwidestar[1,2,5,11,18,25]{(u\odot v)\odot w=(p,s)\odot(t,z)}{\rIE{32,33}}
  \proofstepwidestar[1,2,5,11,18,25]{(p,s)\odot(t,z)=(p,z)}{\FormulaRefAuto{SemigroupCoordinateProductEvaluation}[thm]{1,2,\rAI{31,29}}}
  \proofstepwidestar[1,2,5,11,18,25]{(u\odot v)\odot w=(p,z)}{\rChain{34,35}}
  \proofstepwidestar[11,25]{u\odot w=(p,q)\odot(t,z)}{\rIE{14,28,\rII}}
  \proofstepwidestar[1,2,5,11,25]{(p,q)\odot(t,z)=(p,z)}{\FormulaRefAuto{SemigroupCoordinateProductEvaluation}[thm]{1,2,\rAI{15,29}}}
  \proofstepwidestar[1,2,5,11,25]{u\odot w=(p,z)}{\rChain{37,38}}
  \proofstepwidestar[1,2,5,11,25]{(p,z)=u\odot w}{\FormulaRefAuto{a=b\vdash b=a}[thm]{39}}
  \proofstepwidestar[1,2,5,11,18,25]{(u\odot v)\odot w=u\odot w}{\rIE{40,36}}
  \proofstepwidestar[1,2,5,11,18]{(u\odot v)\odot w=u\odot w}{\rEEStar{24,25,41}}
  \proofstepwidestar[1,2,5,11]{(u\odot v)\odot w=u\odot w}{\rEEStar{17,18,42}}
  \proofstepwidestar[1,2,5]{(u\odot v)\odot w=u\odot w}{\rEEStar{10,11,43}}
\closeproofpartwide
\end{tabproofsplitwide}

\subsection{Rechtecksidentität des Koordinatenprodukts}
\hypertarget{coordinates.proof.SemigroupCoordinateProductRectangularIdentity}{}
\noindent Beweis zu
\FormulaRefAuto{SemigroupCoordinateProductRectangularIdentity}[thm].
\begin{tabproofwide}
  \proofstepwidestar[1]{\SemigroupAlg{A}{\star}}{\rA}
  \proofstepwidestar[2]{e\in A}{\rA}
  \proofstepwidestar[3]{u,v\in K}{\rA}
  \proofstepwidestar[3]{u\in K}{\rAEa{3}}
  \proofstepwidestar[3]{v\in K}{\rAEb{3}}
  \proofstepwidestar[1,2,3]{(u\odot v)\odot u=u\odot u}{\FormulaRefAuto{SemigroupCoordinateProductThreeTermLaw}[thm]{1,2,\rAI{4,\rAI{5,4}}}}
  \proofstepwidestar[1,2,3]{u\odot u=u}{\FormulaRefAuto{SemigroupCoordinateProductIdempotence}[thm]{1,2,4}}
  \proofstepwidestar[1,2,3]{(u\odot v)\odot u=u}{\rChain{6,7}}
\end{tabproofwide}

\Needspace{18\baselineskip}
\section{Die beiden Abbildungen}
\subsection{Typisierung der Koordinatenabbildung}
\hypertarget{coordinates.proof.SemigroupCoordinateMapFunction}{}
\noindent Beweis zu
\FormulaRefAuto{SemigroupCoordinateMapFunction}[thm].
\begin{tabproofwide}
  \proofstepwidestar[1]{\SemigroupAlg{A}{\star}}{\rA}
  \proofstepwidestar[2]{e\in A}{\rA}
  \proofstepwidestar[3]{x\in A}{\rA}
  \proofstepwidestar[1,2,3]{x\star e\in A}{\FormulaRefAuto{SemigroupClosureAxiom}[ax]{1,\rAI{3,2}}}
  \proofstepwidestar[1,2,3]{e\star x\in A}{\FormulaRefAuto{SemigroupClosureAxiom}[ax]{1,\rAI{2,3}}}
  \proofstepwidestar[3]{\exists a\in A\;(x\star e=a\star e)}{\rEI{\rAI{3,\rII}}}
  \proofstepwidestar[3]{\exists b\in A\;(e\star x=e\star b)}{\rEI{\rAI{3,\rII}}}
  \proofstepwidestar[1,2]{\begin{aligned}[t]&(x\star e,e\star x)\in K\leftrightarrow{}\\[-2pt]&\quad(x\star e\in A\land\exists a\in A\;(x\star e=a\star e))\\[-2pt]&\quad\land(e\star x\in A\land\exists b\in A\;(e\star x=e\star b))\end{aligned}}{\FormulaRefAuto{SemigroupCoordinateCarrierMembership}[thm]{1,2}}
  \proofstepwidestar[1,2,3]{(x\star e,e\star x)\in K}{\rRE{\rLREb{8},\rAI{\rAI{4,6},\rAI{5,7}}}}
  \proofstepwidestar[1,2]{\forall x\in A\;((x\star e,e\star x)\in K)}{\rUI{\rRI{3,9}}}
  \proofstepwidestar[1,2]{\chi\colon A\to K}{\FormulaRefAuto{SemigroupCoordinateMapDef}[def]{1,2}}
\end{tabproofwide}

Die vorhergehende Allformel bestätigt den in der Definition angegebenen
Zielbereich. Die Typisierung selbst ist bereits Teil der zitierten
Definition. Beide Komponenten sind durch die binäre Operation bestimmt;
die Zuordnung benötigt keine Wahl von Zeugen.
\subsection{Auswertung der Koordinatenabbildung}
\hypertarget{coordinates.proof.SemigroupCoordinateMapEvaluation}{}
\noindent Beweis zu
\FormulaRefAuto{SemigroupCoordinateMapEvaluation}[thm].
\begin{tabproofwide}
  \proofstepwidestar[1]{\SemigroupAlg{A}{\star}}{\rA}
  \proofstepwidestar[2]{e,x\in A}{\rA}
  \proofstepwidestar[2]{e\in A}{\rAEa{2}}
  \proofstepwidestar[2]{x\in A}{\rAEb{2}}
  \proofstepwidestar[1,2]{\forall x\in A\;(\chi(x)=(x\star e,e\star x))}{\FormulaRefAuto{SemigroupCoordinateMapDef}[def]{1,3}}
  \proofstepwidestar[1,2]{\chi(x)=(x\star e,e\star x)}{\rRE{\rUE{5},4}}
\end{tabproofwide}

\Needspace{14\baselineskip}
\subsection{Typisierung der Multiplikationsabbildung}
\hypertarget{coordinates.proof.SemigroupCoordinateMultiplicationMapFunction}{}
\noindent Beweis zu
\FormulaRefAuto{SemigroupCoordinateMultiplicationMapFunction}[thm].
\begin{tabproofsplitwide}
\proofpartwide{Produktzerlegung und Zielbereich}
  \proofstepwidestar[1]{\SemigroupAlg{A}{\star}}{\rA}
  \proofstepwidestar[2]{e\in A}{\rA}
  \proofstepwidestar[1,2]{K=L\times R}{\FormulaRefAuto{SemigroupCoordinateCarrierDef}[def]{1,2}}
  \proofstepwidestar[1,2]{L\times R=K}{\FormulaRefAuto{a=b\vdash b=a}[thm]{3}}
  \proofstepwidestar[5]{u\in K}{\rA}
  \proofstepwidestar[1,2,5]{u\in L\times R}{\rIE{3,5}}
  \proofstepwidestar[1,2,5]{\exists p\in L\,\exists q\in R\;(u=(p,q))}{\FormulaRefAuto{x\in A\times B\eqvdash\exists a\in A\,\exists b\in B\,x=(a,b)}[thm]{6}}
  \proofstepwidestar[8]{p\in L\land(q\in R\land u=(p,q))}{\rA}
  \proofstepwidestar[8]{p\in L}{\rAEa{8}}
  \proofstepwidestar[8]{q\in R}{\rAEa{\rAEb{8}}}
  \proofstepwidestar[8]{u=(p,q)}{\rAEb{\rAEb{8}}}
  \proofstepwidestar[5,8]{(p,q)\in K}{\rIE{11,5}}
  \proofstepwidestar[1,2]{\begin{aligned}[t](p,q)\in K\leftrightarrow{}&(p\in A\land\exists a\in A\;(p=a\star e))\\[-2pt]&\land(q\in A\land\exists b\in A\;(q=e\star b))\end{aligned}}{\FormulaRefAuto{SemigroupCoordinateCarrierMembership}[thm]{1,2}}
  \proofstepwidestar[1,2,5,8]{\begin{aligned}[t]&(p\in A\land\exists a\in A\;(p=a\star e))\\[-2pt]&\quad\land(q\in A\land\exists b\in A\;(q=e\star b))\end{aligned}}{\rRE{\rLREa{13},12}}
  \proofstepwidestar[1,2,5,8]{p\in A}{\rAEa{\rAEa{14}}}
  \proofstepwidestar[1,2,5,8]{q\in A}{\rAEa{\rAEb{14}}}
  \proofstepwidestar[1,2,5,8]{p\star q\in A}{\FormulaRefAuto{SemigroupClosureAxiom}[ax]{1,\rAI{15,16}}}
  \proofstepwidestar[1,2,5,8]{\mu(p,q)=p\star q}{\FormulaRefAuto{SemigroupCoordinateMultiplicationMapDef}[def]{1,2,12}}
  \proofstepwidestar[8]{\mu(u)=\mu(p,q)}{\rIE{11,\rII}}
  \proofstepwidestar[1,2,5,8]{\mu(u)=p\star q}{\rChain{19,18}}
  \proofstepwidestar[1,2,5,8]{p\star q=\mu(u)}{\FormulaRefAuto{a=b\vdash b=a}[thm]{20}}
  \proofstepwidestar[1,2,5,8]{\mu(u)\in A}{\rIE{21,17}}
  \proofstepwidestar[1,2,5]{\mu(u)\in A}{\rEEStar{7,8,22}}
  \proofstepwidestar[1,2]{\forall u\in K\;(\mu(u)\in A)}{\rUI{\rRI{5,23}}}
  \proofstepwidestar[1,2]{\mu\colon K\to A}{\FormulaRefAuto{SemigroupCoordinateMultiplicationMapDef}[def]{1,2}}
\closeproofpartwide
\end{tabproofsplitwide}

Die Allformel bestätigt den in der Definition angegebenen Zielbereich;
die Typisierung ist bereits Bestandteil dieser Definition.
Die Komponenten eines geordneten Paars sind eindeutig nach
\FormulaRefAuto{(a,b)=(c,d)\eqvdash a=c\land b=d}[thm].
Deshalb hängt \(p\star q\) allein von \(u=(p,q)\) ab, nicht von den
Zeugen für die Mitgliedschaften in \(L\) und \(R\).
\subsection{Auswertung der Multiplikationsabbildung}
\hypertarget{coordinates.proof.SemigroupCoordinateMultiplicationMapEvaluation}{}
\noindent Beweis zu
\FormulaRefAuto{SemigroupCoordinateMultiplicationMapEvaluation}[thm].
\begin{tabproofwide}
  \proofstepwidestar[1]{\SemigroupAlg{A}{\star}}{\rA}
  \proofstepwidestar[2]{e\in A}{\rA}
  \proofstepwidestar[3]{(p,q)\in K}{\rA}
  \proofstepwidestar[1,2,3]{\mu(p,q)=p\star q}{\FormulaRefAuto{SemigroupCoordinateMultiplicationMapDef}[def]{1,2,3}}
\end{tabproofwide}

\Needspace{20\baselineskip}
\section{Isomorphismus und Kriterium}

\subsection{Einordnung als Halbgruppenisomorphismus}
\hypertarget{coordinates.proof.SemigroupCoordinateIsoIsSemigroupIso}{}
\noindent Beweis zu
\FormulaRefAuto{SemigroupCoordinateIsoIsSemigroupIso}[thm].
\begin{tabproofwide}
  \proofstepwidestar[1]{\SemigroupCoordinateIso{A}{\star}{e}}{\rA}
  \proofstepwidestar[1]{\SemigroupAlg{A}{\star}}{\FormulaRefAuto{SemigroupCoordinateIsoSourceAxiom}[ax]{1}}
  \proofstepwidestar[1]{e\in A}{\FormulaRefAuto{SemigroupCoordinateIsoBaseElementAxiom}[ax]{1}}
  \proofstepwidestar[1]{\SemigroupAlg{K}{\odot}}{\FormulaRefAuto{SemigroupCoordinateProductSemigroup}[thm]{2,3}}
  \proofstepwidestar[1]{\chi\colon A\bij K}{\FormulaRefAuto{SemigroupCoordinateIsoBijectiveAxiom}[ax]{1}}
  \proofstepwidestar[1]{\chi\colon A\to K}{\FormulaRefAuto{SemigroupCoordinateMapFunction}[thm]{2,3}}
  \proofstepwidestar[7]{x\in A}{\rA}
  \proofstepwidestar[8]{y\in A}{\rA}
  \proofstepwidestar[1,7,8]{\chi(x\star y)=\chi(x)\odot\chi(y)}{%
    \FormulaRefAuto{SemigroupCoordinateIsoOperationAxiom}[ax]{1,\rAI{7,8}}}
  \proofstepwidestar[1,7]{\forall y\in A\;(\chi(x\star y)=\chi(x)\odot\chi(y))}{\rUI{\rRI{8,9}}}
  \proofstepwidestar[1]{\forall x,y\in A\;(\chi(x\star y)=\chi(x)\odot\chi(y))}{\rUI{\rRI{7,10}}}
  \proofstepwidestar[1]{\SemigroupHom{\chi}{A,\star}{K,\odot}}{\FormulaRefAuto{SemigroupHomDef}[def]{2,4,6,11}}
  \proofstepwidestar[1]{\SemigroupIso{\chi}{A,\star}{K,\odot}}{\FormulaRefAuto{SemigroupIsoDef}[def]{12,5}}
\end{tabproofwide}

\subsection{Idempotenz- und Dreitermkriterium}
\hypertarget{coordinates.proof.SemigroupCoordinateIsoCriterion}{}
\noindent Beweis zu
\FormulaRefAuto{SemigroupCoordinateIsoCriterion}[thm].
Die Rückgewinnung \(\mu(\chi(x))=x\) liefert zunächst die Injektivität.
Für ein beliebiges \(u=(p,q)\in K\) wird anschließend \(p\star q\)
als konkretes Urbild konstruiert. Beide Komponentenzeugen werden vor
Verlassen dieses Teilbeweises entlassen. Die Operationserhaltung folgt
gesondert. Beschränkte Allquantoren werden einzeln spezialisiert und
wieder eingeführt; bei Gleichheitsersetzungen wird die benötigte
Richtung ausdrücklich angegeben.

\Needspace{12\baselineskip}
\begin{tabproofsplitwide}
\proofpartwide{Rückgewinnung: Typisierung und gerichtete Umformungen}
  \proofstepwidestar[1]{\SemigroupAlg{A}{\star}}{%
    \rA}
  \proofstepwidestar[2]{e\in A}{%
    \rA}
  \proofstepwidestar[3]{\forall a\in A\;(a\star a=a)}{%
    \rA}
  \proofstepwidestar[4]{\forall a,b,c\in A\;((a\star b)\star c=a\star c)}{%
    \rA}
  \proofstepwidestar[1,2]{\chi\colon A\to K}{%
    \FormulaRefAuto{SemigroupCoordinateMapFunction}[thm]{1,2}}
  \proofstepwidestar[1,2]{\mu\colon K\to A}{%
    \FormulaRefAuto{SemigroupCoordinateMultiplicationMapFunction}[thm]{1,2}}
  \proofstepwidestar[1,2]{K=L\times R}{%
    \FormulaRefAuto{SemigroupCoordinateCarrierDef}[def]{1,2}}
  \proofstepwidestar[8]{x\in A}{%
    \rA}
  \proofstepwidestar[1,2,8]{x\star e\in A}{%
    \FormulaRefAuto{SemigroupClosureAxiom}[ax]{1,\rAI{8,2}}}
  \proofstepwidestar[1,2,8]{\chi(x)\in K}{%
    \FormulaRefAuto{F\colon A\to B\dsep x\in A\vdash F(x)\in B}[thm]{5,8}}
  \proofstepwidestar[1,2,8]{\chi(x)=(x\star e,e\star x)}{%
    \FormulaRefAuto{SemigroupCoordinateMapEvaluation}[thm]{1,\rAI{2,8}}}
  \proofstepwidestar[1,2,8]{(x\star e,e\star x)\in K}{%
    \rIE{11,10}}
  \proofstepwidestar[1,2,8]{\mu(x\star e,e\star x)=(x\star e)\star(e\star x)}{%
    \FormulaRefAuto{SemigroupCoordinateMultiplicationMapEvaluation}[thm]{1,2,12}}
  \proofstepwidestar[1,2,8]{\mu(\chi(x))=(x\star e)\star(e\star x)}{%
    \rIE{\FormulaRefAuto{a=b\vdash b=a}[thm]{11},13}}
  \proofstepwidestar[1,2,8]{((x\star e)\star e)\star x=(x\star e)\star(e\star x)}{%
    \FormulaRefAuto{SemigroupAssociativityAxiom}[ax]{1,\rAI{9,\rAI{2,8}}}}
  \proofstepwidestar[1,2,8]{(x\star e)\star(e\star x)=((x\star e)\star e)\star x}{%
    \FormulaRefAuto{a=b\vdash b=a}[thm]{15}}
  \proofstepwidestar[2,4,8]{(x\star e)\star e=x\star e}{%
    \rRE{\rUE{\rRE{\rUE{\rRE{\rUE{4},8}},2}},2}}
  \proofstepwidestar[1,2,4,8]{(x\star e)\star(e\star x)=(x\star e)\star x}{%
    \rIE{17,16}}
  \proofstepwidestar[2,4,8]{(x\star e)\star x=x\star x}{%
    \rRE{\rUE{\rRE{\rUE{\rRE{\rUE{4},8}},2}},8}}
  \proofstepwidestar[3,8]{x\star x=x}{%
    \rRE{\rUE{3},8}}
  \proofstepwidestar[1,2,4,8]{\mu(\chi(x))=(x\star e)\star x}{%
    \rIE{18,14}}
  \proofstepwidestar[1,2,4,8]{\mu(\chi(x))=x\star x}{%
    \rIE{19,21}}
  \proofstepwidestar[1,2,3,4,8]{\mu(\chi(x))=x}{%
    \rIE{20,22}}
\nopagebreak
  \proofstepwidestar[1,2,3,4]{\forall x\in A\;(\mu(\chi(x))=x)}{%
    \rUI{\rRI{8,23}}}
\closeproofpartwide
\Needspace{8\baselineskip}
\proofpartwide{Injektivität mit einzeln entlassenen Trägerannahmen}
\setcounter{proofstepnr}{24}
  \proofstepwidestar[25]{x\in A}{%
    \rA}
  \proofstepwidestar[26]{y\in A}{%
    \rA}
  \proofstepwidestar[27]{\chi(x)=\chi(y)}{%
    \rA}
  \proofstepwidestar[1,2,3,4,25]{\mu(\chi(x))=x}{%
    \rRE{\rUE{24},25}}
  \proofstepwidestar[1,2,3,4,26]{\mu(\chi(y))=y}{%
    \rRE{\rUE{24},26}}
  \proofstepwidestar[27]{\mu(\chi(x))=\mu(\chi(y))}{%
    \rIE{27,\rII}}
  \proofstepwidestar[1,2,3,4,25,26,27]{x=y}{%
    \rIE{28,29,30}}
  \proofstepwidestar[1,2,3,4,25,26]{\chi(x)=\chi(y)\rightarrow x=y}{%
    \rRI{27,31}}
  \proofstepwidestar[1,2,3,4,25]{\forall y\in A\;(\chi(x)=\chi(y)\rightarrow x=y)}{%
    \rUI{\rRI{26,32}}}
  \proofstepwidestar[1,2,3,4]{\forall x,y\in A\;(\chi(x)=\chi(y)\rightarrow x=y)}{%
    \rUI{\rRI{25,33}}}
  \proofstepwidestar[1,2,3,4]{\chi\colon A\inj K}{%
    \FormulaRefAuto{Injektive Funktion}[def]{5,34}}
\closeproofpartwide
\Needspace{8\baselineskip}
\proofpartwide{Ein Koordinatenpaar und der Zeuge seiner linken Komponente}
\setcounter{proofstepnr}{35}
  \proofstepwidestar[36]{u\in K}{%
    \rA}
  \proofstepwidestar[1,2,36]{u\in L\times R}{%
    \rIE{7,36}}
  \proofstepwidestar[1,2,36]{\exists p\in L\,\exists q\in R\;(u=(p,q))}{%
    \FormulaRefAuto{x\in A\times B\eqvdash\exists a\in A\,\exists b\in B\,x=(a,b)}[thm]{37}}
  \proofstepwidestar[39]{p\in L\land(q\in R\land u=(p,q))}{%
    \rA}
  \proofstepwidestar[39]{u=(p,q)}{%
    \rAEb{\rAEb{39}}}
  \proofstepwidestar[36,39]{(p,q)\in K}{%
    \rIE{40,36}}
  \proofstepwidestar[1,2]{\begin{aligned}[t]&(p,q)\in K\leftrightarrow\\[-2pt]&\quad(p\in A\land\exists a\in A\;(p=a\star e))\\[-2pt]&\qquad\land(q\in A\land\exists b\in A\;(q=e\star b))\end{aligned}}{%
    \FormulaRefAuto{SemigroupCoordinateCarrierMembership}[thm]{1,2}}
  \proofstepwidestar[1,2,36,39]{\begin{aligned}[t]&(p\in A\land\exists a\in A\;(p=a\star e))\\[-2pt]&\quad\land(q\in A\land\exists b\in A\;(q=e\star b))\end{aligned}}{%
    \rRE{\rLREa{42},41}}
  \proofstepwidestar[1,2,36,39]{p\in A}{%
    \rAEa{\rAEa{43}}}
  \proofstepwidestar[1,2,36,39]{q\in A}{%
    \rAEa{\rAEb{43}}}
  \proofstepwidestar[1,2,36,39]{\exists a\in A\;(p=a\star e)}{%
    \rAEb{\rAEa{43}}}
  \proofstepwidestar[47]{a\in A\land p=a\star e}{%
    \rA}
  \proofstepwidestar[47]{a\in A}{%
    \rAEa{47}}
  \proofstepwidestar[47]{p=a\star e}{%
    \rAEb{47}}
  \proofstepwidestar[2,4,47]{(a\star e)\star e=a\star e}{%
    \rRE{\rUE{\rRE{\rUE{\rRE{\rUE{4},48}},2}},2}}
  \proofstepwidestar[47]{a\star e=p}{%
    \FormulaRefAuto{a=b\vdash b=a}[thm]{49}}
  \proofstepwidestar[2,4,47]{p\star e=p}{%
    \rIE{51,50}}
  \proofstepwidestar[1,2,4,36,39]{p\star e=p}{%
    \rEE{46,47,52}}
\closeproofpartwide
\Needspace{8\baselineskip}
\proofpartwide{Der Zeuge der rechten Komponente}
\setcounter{proofstepnr}{53}
  \proofstepwidestar[1,2,36,39]{\exists b\in A\;(q=e\star b)}{%
    \rAEb{\rAEb{43}}}
  \proofstepwidestar[55]{b\in A\land q=e\star b}{%
    \rA}
  \proofstepwidestar[55]{b\in A}{%
    \rAEa{55}}
  \proofstepwidestar[55]{q=e\star b}{%
    \rAEb{55}}
  \proofstepwidestar[2,3]{e\star e=e}{%
    \rRE{\rUE{3},2}}
  \proofstepwidestar[1,2,55]{(e\star e)\star b=e\star(e\star b)}{%
    \FormulaRefAuto{SemigroupAssociativityAxiom}[ax]{1,\rAI{2,\rAI{2,56}}}}
  \proofstepwidestar[1,2,55]{e\star(e\star b)=(e\star e)\star b}{%
    \FormulaRefAuto{a=b\vdash b=a}[thm]{59}}
  \proofstepwidestar[1,2,55]{e\star q=(e\star e)\star b}{%
    \rIE{\FormulaRefAuto{a=b\vdash b=a}[thm]{57},60}}
  \proofstepwidestar[1,2,3,55]{e\star q=e\star b}{%
    \rIE{58,61}}
  \proofstepwidestar[1,2,3,55]{e\star q=q}{%
    \rIE{\FormulaRefAuto{a=b\vdash b=a}[thm]{57},62}}
  \proofstepwidestar[1,2,3,36,39]{e\star q=q}{%
    \rEE{54,55,63}}
\closeproofpartwide
\Needspace{8\baselineskip}
\proofpartwide{Das Produkt der Komponenten ist ein Urbild}
\setcounter{proofstepnr}{64}
  \proofstepwidestar[1,2,36,39]{p\star q\in A}{%
    \FormulaRefAuto{SemigroupClosureAxiom}[ax]{1,\rAI{44,45}}}
  \proofstepwidestar[1,2,4,36,39]{(p\star q)\star e=p\star e}{%
    \rRE{\rUE{\rRE{\rUE{\rRE{\rUE{4},44}},45}},2}}
  \proofstepwidestar[1,2,4,36,39]{(p\star q)\star e=p}{%
    \rIE{53,66}}
  \proofstepwidestar[1,2,36,39]{(e\star p)\star q=e\star(p\star q)}{%
    \FormulaRefAuto{SemigroupAssociativityAxiom}[ax]{1,\rAI{2,\rAI{44,45}}}}
  \proofstepwidestar[1,2,36,39]{e\star(p\star q)=(e\star p)\star q}{%
    \FormulaRefAuto{a=b\vdash b=a}[thm]{68}}
  \proofstepwidestar[1,2,4,36,39]{(e\star p)\star q=e\star q}{%
    \rRE{\rUE{\rRE{\rUE{\rRE{\rUE{4},2}},44}},45}}
  \proofstepwidestar[1,2,3,4,36,39]{e\star(p\star q)=q}{%
    \rIE{70,64,69}}
  \proofstepwidestar[1,2,36,39]{\chi(p\star q)=((p\star q)\star e,e\star(p\star q))}{%
    \FormulaRefAuto{SemigroupCoordinateMapEvaluation}[thm]{1,\rAI{2,65}}}
  \proofstepwidestar[1,2,3,4,36,39]{\chi(p\star q)=(p,q)}{%
    \rIE{67,71,72}}
  \proofstepwidestar[1,2,3,4,36,39]{\chi(p\star q)=u}{%
    \rIE{\FormulaRefAuto{a=b\vdash b=a}[thm]{40},73}}
  \proofstepwidestar[1,2,3,4,36,39]{\exists x\in A\;(\chi(x)=u)}{%
    \rEI{\rAI{65,74}}}
  \proofstepwidestar[1,2,3,4,36]{\exists x\in A\;(\chi(x)=u)}{%
    \rEEStar{38,39,75}}
  \proofstepwidestar[1,2,3,4]{\forall u\in K\,\exists x\in A\;(\chi(x)=u)}{%
    \rUI{\rRI{36,76}}}
\nopagebreak
  \proofstepwidestar[1,2,3,4]{\chi\colon A\sur K}{%
    \FormulaRefAuto{Surjektive Funktion}[def]{5,77}}
\nopagebreak
  \proofstepwidestar[1,2,3,4]{\chi\colon A\bij K}{%
    \FormulaRefAuto{Bijektive Funktion}[def]{5,34,77}}
\closeproofpartwide
\Needspace{8\baselineskip}
\proofpartwide{Operationserhaltung: Träger und beide Koordinaten}
\setcounter{proofstepnr}{79}
  \proofstepwidestar[80]{x\in A}{%
    \rA}
  \proofstepwidestar[81]{y\in A}{%
    \rA}
  \proofstepwidestar[1,80,81]{x\star y\in A}{%
    \FormulaRefAuto{SemigroupClosureAxiom}[ax]{1,\rAI{80,81}}}
  \proofstepwidestar[1,2,80]{\chi(x)\in K}{%
    \FormulaRefAuto{F\colon A\to B\dsep x\in A\vdash F(x)\in B}[thm]{5,80}}
  \proofstepwidestar[1,2,81]{\chi(y)\in K}{%
    \FormulaRefAuto{F\colon A\to B\dsep x\in A\vdash F(x)\in B}[thm]{5,81}}
  \proofstepwidestar[1,2,80]{\chi(x)=(x\star e,e\star x)}{%
    \FormulaRefAuto{SemigroupCoordinateMapEvaluation}[thm]{1,\rAI{2,80}}}
  \proofstepwidestar[1,2,81]{\chi(y)=(y\star e,e\star y)}{%
    \FormulaRefAuto{SemigroupCoordinateMapEvaluation}[thm]{1,\rAI{2,81}}}
  \proofstepwidestar[1,2,80]{(x\star e,e\star x)\in K}{%
    \rIE{85,83}}
  \proofstepwidestar[1,2,81]{(y\star e,e\star y)\in K}{%
    \rIE{86,84}}
  \proofstepwidestar[2,4,80,81]{(x\star y)\star e=x\star e}{%
    \rRE{\rUE{\rRE{\rUE{\rRE{\rUE{4},80}},81}},2}}
  \proofstepwidestar[1,2,80,81]{(e\star x)\star y=e\star(x\star y)}{%
    \FormulaRefAuto{SemigroupAssociativityAxiom}[ax]{1,\rAI{2,\rAI{80,81}}}}
  \proofstepwidestar[1,2,80,81]{e\star(x\star y)=(e\star x)\star y}{%
    \FormulaRefAuto{a=b\vdash b=a}[thm]{90}}
  \proofstepwidestar[2,4,80,81]{(e\star x)\star y=e\star y}{%
    \rRE{\rUE{\rRE{\rUE{\rRE{\rUE{4},2}},80}},81}}
  \proofstepwidestar[1,2,4,80,81]{e\star(x\star y)=e\star y}{%
    \rIE{92,91}}
\closeproofpartwide
\Needspace{8\baselineskip}
\proofpartwide{Operationserhaltung und Schluss}
\setcounter{proofstepnr}{93}
  \proofstepwidestar[1,2,80,81]{\chi(x\star y)=((x\star y)\star e,e\star(x\star y))}{%
    \FormulaRefAuto{SemigroupCoordinateMapEvaluation}[thm]{1,\rAI{2,82}}}
  \proofstepwidestar[1,2,4,80,81]{\chi(x\star y)=(x\star e,e\star y)}{%
    \rIE{89,93,94}}
  \proofstepwidestar[1,2,80,81]{\begin{aligned}[t]&(x\star e,e\star x)\odot(y\star e,e\star y)\\[-2pt]&\quad=(x\star e,e\star y)\end{aligned}}{%
    \FormulaRefAuto{SemigroupCoordinateProductEvaluation}[thm]{1,2,\rAI{87,88}}}
  \proofstepwidestar[1,2,80,81]{\chi(x)\odot\chi(y)=(x\star e,e\star y)}{%
    \rIE{\FormulaRefAuto{a=b\vdash b=a}[thm]{85},\FormulaRefAuto{a=b\vdash b=a}[thm]{86},96}}
  \proofstepwidestar[1,2,80,81]{(x\star e,e\star y)=\chi(x)\odot\chi(y)}{%
    \FormulaRefAuto{a=b\vdash b=a}[thm]{97}}
  \proofstepwidestar[1,2,4,80,81]{\chi(x\star y)=\chi(x)\odot\chi(y)}{%
    \rIE{98,95}}
\nopagebreak
  \proofstepwidestar[1,2,4,80]{\forall y\in A\;(\chi(x\star y)=\chi(x)\odot\chi(y))}{%
    \rUI{\rRI{81,99}}}
\nopagebreak
  \proofstepwidestar[1,2,4]{\forall x,y\in A\;(\chi(x\star y)=\chi(x)\odot\chi(y))}{%
    \rUI{\rRI{80,100}}}
\nopagebreak
  \proofstepwidestar[1,2,3,4]{\SemigroupCoordinateIso{A}{\star}{e}}{%
    \FormulaRefAuto{SemigroupCoordinateIsoDef}[def]{1,2,79,101}}
\closeproofpartwide
\end{tabproofsplitwide}

\Needspace{28\baselineskip}
\section{Rücktransport der drei Identitäten}

\subsection{Idempotenz der Ausgangshalbgruppe}
\hypertarget{coordinates.proof.SemigroupCoordinateIsoIdempotence}{}
\noindent Beweis zu
\FormulaRefAuto{SemigroupCoordinateIsoIdempotence}[thm].
\begin{tabproofwide}
  \proofstepwidestar[1]{\SemigroupCoordinateIso{A}{\star}{e}}{\rA}
  \proofstepwidestar[1]{\SemigroupAlg{A}{\star}}{\FormulaRefAuto{SemigroupCoordinateIsoSourceAxiom}[ax]{1}}
  \proofstepwidestar[1]{e\in A}{\FormulaRefAuto{SemigroupCoordinateIsoBaseElementAxiom}[ax]{1}}
  \proofstepwidestar[1]{\chi\colon A\bij K}{\FormulaRefAuto{SemigroupCoordinateIsoBijectiveAxiom}[ax]{1}}
  \proofstepwidestar[1]{\chi\colon A\inj K}{\FormulaRefAuto{F\colon A\bij B \vdash F\colon A\inj B}[thm]{4}}
  \proofstepwidestar[1]{\chi\colon A\to K}{\FormulaRefAuto{SemigroupCoordinateMapFunction}[thm]{2,3}}
  \proofstepwidestar[7]{x\in A}{\rA}
  \proofstepwidestar[1,7]{x\star x\in A}{\FormulaRefAuto{SemigroupClosureAxiom}[ax]{2,\rAI{7,7}}}
  \proofstepwidestar[1,7]{\chi(x)\in K}{\FormulaRefAuto{F\colon A\to B\dsep x\in A\vdash F(x)\in B}[thm]{6,7}}
  \proofstepwidestar[1,7]{\chi(x\star x)=\chi(x)\odot\chi(x)}{\FormulaRefAuto{SemigroupCoordinateIsoOperationAxiom}[ax]{1,\rAI{7,7}}}
  \proofstepwidestar[1,7]{\chi(x)\odot\chi(x)=\chi(x)}{\FormulaRefAuto{SemigroupCoordinateProductIdempotence}[thm]{2,3,9}}
  \proofstepwidestar[1,7]{\chi(x\star x)=\chi(x)}{\rChain{10,11}}
  \proofstepwidestar[1,7]{x\star x=x}{\FormulaRefAuto{Injektivität}[ax]{5,8,7,12}}
  \proofstepwidestar[1]{\forall x\in A\;(x\star x=x)}{\rUI{\rRI{7,13}}}
\end{tabproofwide}

\Needspace{12\baselineskip}
\subsection{Dreitermgesetz der Ausgangshalbgruppe}
\hypertarget{coordinates.proof.SemigroupCoordinateIsoThreeTermLaw}{}
\noindent Beweis zu
\FormulaRefAuto{SemigroupCoordinateIsoThreeTermLaw}[thm].
\begin{tabproofsplitwide}
\proofpartwide{Typisierung und Operationserhaltung}
  \proofstepwidestar[1]{\SemigroupCoordinateIso{A}{\star}{e}}{\rA}
  \proofstepwidestar[1]{\SemigroupAlg{A}{\star}}{\FormulaRefAuto{SemigroupCoordinateIsoSourceAxiom}[ax]{1}}
  \proofstepwidestar[1]{e\in A}{\FormulaRefAuto{SemigroupCoordinateIsoBaseElementAxiom}[ax]{1}}
  \proofstepwidestar[1]{\chi\colon A\bij K}{\FormulaRefAuto{SemigroupCoordinateIsoBijectiveAxiom}[ax]{1}}
  \proofstepwidestar[1]{\chi\colon A\inj K}{\FormulaRefAuto{F\colon A\bij B \vdash F\colon A\inj B}[thm]{4}}
  \proofstepwidestar[1]{\chi\colon A\to K}{\FormulaRefAuto{SemigroupCoordinateMapFunction}[thm]{2,3}}
  \proofstepwidestar[7]{x\in A}{\rA}
  \proofstepwidestar[8]{y\in A}{\rA}
  \proofstepwidestar[9]{z\in A}{\rA}
  \proofstepwidestar[1,7,8]{x\star y\in A}{\FormulaRefAuto{SemigroupClosureAxiom}[ax]{2,\rAI{7,8}}}
  \proofstepwidestar[1,7,8,9]{(x\star y)\star z\in A}{\FormulaRefAuto{SemigroupClosureAxiom}[ax]{2,\rAI{10,9}}}
  \proofstepwidestar[1,7,9]{x\star z\in A}{\FormulaRefAuto{SemigroupClosureAxiom}[ax]{2,\rAI{7,9}}}
  \proofstepwidestar[1,7]{\chi(x)\in K}{\FormulaRefAuto{F\colon A\to B\dsep x\in A\vdash F(x)\in B}[thm]{6,7}}
  \proofstepwidestar[1,8]{\chi(y)\in K}{\FormulaRefAuto{F\colon A\to B\dsep x\in A\vdash F(x)\in B}[thm]{6,8}}
  \proofstepwidestar[1,9]{\chi(z)\in K}{\FormulaRefAuto{F\colon A\to B\dsep x\in A\vdash F(x)\in B}[thm]{6,9}}
  \proofstepwidestar[1,7,8]{\chi(x\star y)=\chi(x)\odot\chi(y)}{\FormulaRefAuto{SemigroupCoordinateIsoOperationAxiom}[ax]{1,\rAI{7,8}}}
  \proofstepwidestar[1,7,9]{\chi(x\star z)=\chi(x)\odot\chi(z)}{\FormulaRefAuto{SemigroupCoordinateIsoOperationAxiom}[ax]{1,\rAI{7,9}}}
\closeproofpartwide
\Needspace{8\baselineskip}
\proofpartwide{Gleichheit der Bilder und Injektivität}
\setcounter{proofstepnr}{17}
  \proofstepwide[1,7,8,9]{\chi((x\star y)\star z)}{=}{\chi(x\star y)\odot\chi(z)}{%
    \FormulaRefAuto{SemigroupCoordinateIsoOperationAxiom}[ax]{1,\rAI{10,9}}}
  \proofstepwide[1,7,8]{}{=}{(\chi(x)\odot\chi(y))\odot\chi(z)}{\rIE{16,\rII}}
  \proofstepwide[1,7,8,9]{}{=}{\chi(x)\odot\chi(z)}{\FormulaRefAuto{SemigroupCoordinateProductThreeTermLaw}[thm]{2,3,\rAI{13,\rAI{14,15}}}}
  \proofstepwide[1,7,9]{}{=}{\chi(x\star z)}{\rIE{17,\rII}}
  \proofstepwidestar[1,7,8,9]{\chi((x\star y)\star z)=\chi(x\star z)}{\rChain{18,19,20,21}}
  \proofstepwidestar[1,7,8,9]{(x\star y)\star z=x\star z}{\FormulaRefAuto{Injektivität}[ax]{5,11,12,22}}
  \proofstepwidestar[1,7,8]{\forall z\in A\;((x\star y)\star z=x\star z)}{\rUI{\rRI{9,23}}}
  \proofstepwidestar[1,7]{\forall y,z\in A\;((x\star y)\star z=x\star z)}{\rUI{\rRI{8,24}}}
  \proofstepwidestar[1]{\forall x,y,z\in A\;((x\star y)\star z=x\star z)}{\rUI{\rRI{7,25}}}
\closeproofpartwide
\end{tabproofsplitwide}

\Needspace{24\baselineskip}
\subsection{Rechtecksidentität der Ausgangshalbgruppe}
\hypertarget{coordinates.proof.SemigroupCoordinateIsoRectangularIdentity}{}
\noindent Beweis zu
\FormulaRefAuto{SemigroupCoordinateIsoRectangularIdentity}[thm].
\begin{tabproofwide}
  \proofstepwidestar[1]{\SemigroupCoordinateIso{A}{\star}{e}}{\rA}
  \proofstepwidestar[1]{\forall a\in A\;(a\star a=a)}{\FormulaRefAuto{SemigroupCoordinateIsoIdempotence}[thm]{1}}
  \proofstepwidestar[1]{\forall a,b,c\in A\;((a\star b)\star c=a\star c)}{\FormulaRefAuto{SemigroupCoordinateIsoThreeTermLaw}[thm]{1}}
  \proofstepwidestar[4]{x\in A}{\rA}
  \proofstepwidestar[5]{y\in A}{\rA}
  \proofstepwidestar[1,4]{\forall b,c\in A\;((x\star b)\star c=x\star c)}{\rRE{4,\rUE{3}}}
  \proofstepwidestar[1,4,5]{\forall c\in A\;((x\star y)\star c=x\star c)}{\rRE{5,\rUE{6}}}
  \proofstepwidestar[1,4,5]{(x\star y)\star x=x\star x}{\rRE{4,\rUE{7}}}
  \proofstepwidestar[1,4]{x\star x=x}{\rRE{4,\rUE{2}}}
  \proofstepwidestar[1,4,5]{(x\star y)\star x=x}{\rChain{8,9}}
  \proofstepwidestar[1,4]{\forall y\in A\;((x\star y)\star x=x)}{\rUI{\rRI{5,10}}}
  \proofstepwidestar[1]{\forall x,y\in A\;((x\star y)\star x=x)}{\rUI{\rRI{4,11}}}
\end{tabproofwide}
